\section{An Execution Substrate for Mixture-of-Models}
\label{sec:reference-engine}

\vsr{}\footnote{Open-source implementation:
\url{https://github.com/vllm-project/semantic-router}.} provides an initial
execution substrate for MoM. Current releases expose selection and bounded
multi-model policies through standard model APIs, driven by typed request
analysis and declared model references, with evaluation and replay hooks around
the resulting decision. This demonstrates that heterogeneous model identities
and multi-call topologies can execute behind one model invocation rather than
inside an agent harness. The portable bundle, environment-owned binding and
lock, topology-neutral interpreter, and canonical event DAG proposed in
Section~\ref{sec:artifact} remain lifecycle extensions rather than shipped
features.

\begin{figure}[htbp]
  \centering
  \resizebox{0.98\linewidth}{!}{\begin{tikzpicture}[
  x=1cm,y=1cm,font=\sffamily,>=Latex,
  flow/.style={-{Latex[length=2.0mm]},draw=MoMGray,line width=0.72pt},
  exchange/.style={<->,>=Latex,draw=MoMGray,line width=0.68pt},
  boundary/.style={draw=MoMBorder,fill=MoMBlueLight!13,rounded corners=5pt,
    line width=0.9pt},
  box/.style={draw=MoMBorder,fill=white,rounded corners=2.8pt,
    line width=0.64pt,minimum height=0.86cm,inner xsep=6pt,
    inner ysep=3pt,align=center,text=MoMText},
  dark/.style={box,draw=none,fill=MoMDark,text=white},
  identity/.style={box,draw=none,fill=MoMBlueLight},
  response/.style={box,draw=MoMDark,fill=MoMPurple,line width=0.76pt},
  title/.style={font=\sffamily\bfseries\fontsize{9.5}{10.3}\selectfont,
    text=MoMText},
  meta/.style={font=\sffamily\fontsize{6.5}{7.2}\selectfont,text=MoMMuted}
]

\newcommand{\momfigmain}{\bfseries\fontsize{8.2}{9.0}\selectfont}
\newcommand{\momfigsub}{\fontsize{6.7}{7.4}\selectfont\color{MoMMuted}}
\newcommand{\momfiglightsub}{\fontsize{6.7}{7.4}\selectfont\color{white}}

% One proposition: a standard model call crosses a single deployed MoM
% boundary containing both the lifecycle engine and its bound constituents.
\draw[boundary] (3.02,0.82) rectangle (12.78,4.10);
\node[title] at (7.90,3.82) {Target deployed MoM boundary};
\node[meta] at (7.90,3.48)
  {versioned identity + behavior variant \;\textbullet\; binding + lock};

\node[identity,minimum width=2.35cm,minimum height=0.96cm] (request) at (1.38,2.45)
  {{\bfseries\fontsize{8.0}{8.8}\selectfont Standard call}\\[-1pt]
   {\fontsize{6.3}{7.0}\selectfont\ttfamily model=mom-v1-ultra}};

\node[dark,minimum width=4.35cm,minimum height=1.04cm] (engine) at (7.90,2.55)
  {{\momfigmain vLLM Semantic Router}\\[-1pt]
   {\momfiglightsub admit \;\textbullet\; decide \;\textbullet\; execute \;\textbullet\; validate}};

\node[response,minimum width=1.75cm,minimum height=0.90cm] (answer) at (14.02,2.45)
  {{\momfigmain One model}\\[-1pt]{\momfigsub response}};

\draw[flow] (request) -- (engine);
\draw[flow] (engine) -- (answer);

\node[box,fill=MoMBlueLight,draw=MoMBorder,minimum width=4.10cm,
  minimum height=0.80cm] (pool) at (7.90,1.42)
  {{\momfigmain Bound constituents}\\[-1pt]
   {\momfigsub local \;\textbullet\; private \;\textbullet\; cloud}};

\draw[exchange] (engine.south) -- (pool.north);

\end{tikzpicture}
}
  \caption{\textbf{The model-invocation boundary of a deployed MoM.} A caller
  issues one standard model request and receives one response, while the engine
  admits, composes, and observes bounded calls to bound constituents. Current
  provider mappings resolve configured names to backend targets. Portable
  identity and environment-specific resolution complete the proposed lifecycle
  boundary.}
  \label{fig:reference-engine}
\end{figure}

\subsection{Why the control point is below agents}

An agent runtime owns application semantics: goals, task decomposition,
interaction state, and the intent to invoke a model. A MoM engine owns a
different contract: whether model-level computation is admissible, how it is
resolved and scheduled, which finite budget it may consume, and how the result
is attributed. The model-invocation boundary is the narrowest common seam at
which this contract can cover direct inference, retrieval-augmented generation,
agent loops, batch jobs, and embedded systems without depending on one agent
framework.

At this seam, a complete engine can enforce authorization and residency,
resolve logical references against deployment evidence, reserve budgets, choose
legal fallbacks, and account for every attempt. Current \vsr{} implements only
part of that contract; the full admissibility and attribution boundary belongs
to the proposed MoM format. When deployed as the exclusive model ingress,
this placement prevents direct calls from bypassing policy and avoids
reproducing governance logic in every agent harness.

The boundary remains closed-world: one invocation must use only declared models
and operators, consume finite resources, and return one terminal response.
An agent may invoke a MoM or request a bounded workflow, but persistent goals,
task queues, and application state remain outside the logical model. Agents are
clients of the model infrastructure, not its portability or accounting
boundary.

Figure~\ref{fig:deployment-envelopes} shows why this placement matters. One
logical MoM can retain its interface, graph, preferences, and constraints while
different bindings realize it on a robot, at an edge site, in a sovereign data
center, or through managed cloud services.

\begin{figure}[htbp]
  \centering
  \resizebox{\linewidth}{!}{\begin{tikzpicture}[
  x=1cm,
  y=1cm,
  >=Latex,
  artifact/.style={
    draw=MoMDark,
    fill=MoMDark,
    text=white,
    rounded corners=2pt,
    minimum height=0.82cm,
    text width=15.25cm,
    align=center,
    font=\sffamily\bfseries\small
  },
  panel/.style={
    draw=MoMBorder,
    fill=white,
    rounded corners=3pt,
    minimum width=3.55cm,
    minimum height=3.25cm
  },
  paneltitle/.style={
    draw=none,
    fill=MoMBlueLight,
    rounded corners=2pt,
    minimum width=3.17cm,
    minimum height=0.48cm,
    align=center,
    font=\sffamily\bfseries\scriptsize,
    text=MoMText
  },
  component/.style={
    draw=MoMBorder,
    fill=white,
    rounded corners=1.5pt,
    minimum width=2.82cm,
    minimum height=0.42cm,
    align=center,
    font=\sffamily\scriptsize,
    text=MoMText
  },
  active/.style={
    component,
    draw=MoMDark,
    fill=MoMBlue
  },
  note/.style={
    align=center,
    font=\sffamily\fontsize{6.1}{6.8}\selectfont,
    text=MoMMuted
  },
  bind/.style={
    -{Latex[length=2.2mm]},
    draw=MoMGray,
    line width=0.65pt
  }
]

\node[artifact] (artifact) at (8.0,5.05)
  {Versioned MoM artifact\\[-1pt]
   \mdseries\scriptsize typed graph $\cdot$ behavior variant $\cdot$ request-preference domain
   $\cdot$ contracts $\cdot$ bounds $\cdot$ content identity};

\foreach \x in {2.0,6.0,10.0,14.0} {
  \draw[bind] (\x,4.63) -- node[right=1pt,note] {bind + lock} (\x,4.12);
}

\node[panel] at (2.0,2.47) {};
\node[paneltitle] at (2.0,3.72) {Robot / on-device};
\node[component] at (2.0,3.02) {sensor + local context};
\node[active] at (2.0,2.42) {compact multimodal model};
\node[component] at (2.0,1.82) {NPU $\cdot$ CPU};
\node[note] at (2.0,1.22) {offline / low latency / bounded};

\node[panel] at (6.0,2.47) {};
\node[paneltitle] at (6.0,3.72) {Edge site};
\node[component] at (6.0,3.02) {site-private data};
\node[active] at (6.0,2.42) {domain specialist};
\node[component] at (6.0,1.82) {NPU $\cdot$ GPU};
\node[note] at (6.0,1.22) {private / resilient / intermittent};

\node[panel] at (10.0,2.47) {};
\node[paneltitle,fill=MoMPurple!32] at (10.0,3.72) {Sovereign data center};
\node[component] at (10.0,3.02) {private data + retrieval};
\node[active,fill=MoMPurple!55] at (10.0,2.42) {open + domain models};
\node[component] at (10.0,1.82) {GPU cluster};
\node[note] at (10.0,1.22) {jurisdiction / audit / control};

\node[panel] at (14.0,2.47) {};
\node[paneltitle] at (14.0,3.72) {Public cloud};
\node[component] at (14.0,3.02) {managed frontier models};
\node[active] at (14.0,2.42) {elastic model services};
\node[component] at (14.0,1.82) {API $\cdot$ cloud GPU};
\node[note] at (14.0,1.22) {elastic / managed / frontier};

\draw[draw=MoMBorder,line width=0.6pt] (0.22,0.76) -- (15.78,0.76);
\node[note,anchor=west] at (0.22,0.50)
  {tighter locality and physical control};
\node[note,anchor=east] at (15.78,0.50)
  {greater elasticity and managed capability};
\node[font=\sffamily\bfseries\scriptsize,text=MoMDark] at (8.0,0.50)
  {one model interface $\cdot$ explicit admissibility $\cdot$ attributable trace};

\end{tikzpicture}
}
  \caption{\textbf{One logical MoM, multiple target deployment envelopes.}
  Environment-owned bindings realize the same versioned specification using
  on-device models, an edge specialist, a jurisdiction-bound private pool, or
  managed cloud capabilities. Portability preserves logical and control
  semantics, not identical outputs; every substitution and measured outcome
  remains explicit behind one model interface.}
  \label{fig:deployment-envelopes}
\end{figure}

\FloatBarrier
\subsection{From routing engine to composite-model lifecycle}

The implementation already separates routing-owned semantics from provider
access details. Configured model references participate in decisions and
bounded multi-call policies, while provider entries map those names to concrete
backends. This is sufficient to execute declared composites, but not for
the full MoM admissibility contract. The proposed binder adds an
environment-owned, fail-closed check that every realized call satisfies
deployment evidence and hard constraints before each call is issued. The
change is a lifecycle boundary, not a claim that a larger
configuration alone constitutes a model.

The distinction is concrete for accelerators. Current platform selection
governs the router's own execution image---CPU by default, or an AMD/ROCm or
NVIDIA/CUDA variant---rather than provisioning constituent LLM endpoints. In
the proposed lifecycle, accelerator choice is resolved per target: one trace
may combine a ROCm-served local model, a CUDA specialist, a CPU verifier, and a
managed provider. Wire transport, API adapter, runtime, leaf artifact,
accelerator, and deployment driver are independent compatibility dimensions.
Existing endpoints are the first binding driver; managed deployment follows
behind the same resolver interface.

The user-facing consequence is intentionally simple. Current releases can
expose configured routed and multi-call policies through ordinary model slugs
such as
\texttt{vllm-sr/auto}, \texttt{vllm-sr/fusion},
\texttt{vllm-sr/remom}, and \texttt{vllm-sr/flow}. The coordinates in
Table~\ref{tab:profile-identities} illustrate the next step: a registry entry
would pin the entire logical MoM and its behavior variant, just as a
model version pins a conventional serving target.

\begin{table}[H]
\caption{Illustrative versioned MoM identities exposed through the standard
model field.}
\label{tab:profile-identities}
\centering
\footnotesize
\setlength{\tabcolsep}{4pt}
\begin{tabular}{L{0.43\linewidth}L{0.20\linewidth}L{0.27\linewidth}}
\toprule
Model coordinate & Behavior variant & Default objective or guard \\
\midrule
\texttt{vllm-sr/mom-v1-flash:1.0.0} & Speed-first & Minimize expected latency \\
\texttt{vllm-sr/mom-v1-light:1.0.0} & Cost-first & Minimize expected cost above a quality floor \\
\texttt{vllm-sr/mom-v1-ultra:1.0.0} & Accuracy-first & Maximize utility under a declared budget \\
\texttt{vllm-sr/mom-v1-halu:1.0.0} & Grounding-first & Require evidence checks and fail-closed fallback \\
\texttt{vllm-sr/mom-v1-secu:1.0.0} & Security-first & Require jailbreak and PII checks \\
\bottomrule
\end{tabular}
\end{table}

To a client, these identities remain values of the standard \texttt{model}
field and return one response. Internally, their policies may select, escalate,
deliberate in parallel, or execute a validated bounded workflow. The project
also supplies the lifecycle hooks needed to study this behavior as one object:
control policies can be trained without gradients through constituents,
complete paths can be evaluated through the serving interface, replay associates
requests with decisions and outcomes, and an offline planner converts routed
demand into capacity alternatives. These surfaces are currently connected
operationally rather than by one content identity. The proposed bundle, lock,
and event schema make that identity the unit of promotion, evaluation, replay,
and cross-engine conformance.
